\documentclass[../../main.tex]{subfiles}% 注意这里的文件路径不能用 ./main.tex ,否则用latexmk编译子文件会报错
\graphicspath{{\subfix{./image/}}{\subfix{../../image/}}} % 指定图片引用路径,后续可以直接使用图片文件名
% 如果图片存放在上级目录,则使用 ../../image/ 作为备用路径

% 例如:
% \begin{figure}[H]
% \centering
% \includegraphics[width=0.3\textwidth]{图.png}
% \caption{}
% \label{figure:图}
% \end{figure}
% 注意:上述\label{}一定要放在\caption{}之后,否则引用图片序号会只会显示??.

\begin{document}

\section{实例}

\begin{examples}[较著名的一些线性偏微分方程]
\begin{enumerate}[(1)]
\item Laplace 方程
\begin{align*}
\Delta u = 0;
\end{align*}

\item 特征值方程
\begin{align*}
\Delta u + \lambda u = 0 \quad (\lambda \text{ 为常数});
\end{align*}

\item 热方程
\begin{align*}
u_t - a^2 \Delta u = 0 \quad (a > 0 \text{ 为常数});
\end{align*}

\item Schr\"{o}dinger 方程
\begin{align*}
u_t - i \Delta u = 0;
\end{align*}

\item Kolmogorov 方程
\begin{align*}
u_t - \sum_{i,j=1}^{n} a_{ij} u_{x_i x_j} + \sum_{i=1}^{n} b_i u_{x_i} = 0,
\end{align*}
其中 $a_{ij}$, $b_i$ $(i,j=1,2,\cdots,n)$ 为常数；

\item Fokker-Planck 方程
\begin{align*}
u_t - \sum_{i,j=1}^{n} (a_{ij} u)_{x_i x_j} + \sum_{i=1}^{n} (b_i u)_{x_i} = 0,
\end{align*}
其中 $a_{ij}$, $b_i$ $(i,j=1,2,\cdots,n)$ 为常数；

\item 输运方程
\begin{align*}
u_t + \sum_{i=1}^{n} b_i u_{x_i} = 0,
\end{align*}
其中 $b_i$ $(i=1,2,\cdots,n)$ 为常数；

\item 波动方程
\begin{align*}
u_{tt} - a^2 \Delta u = 0 \quad (a > 0 \text{ 为常数});
\end{align*}

\item 电报方程
\begin{align*}
u_{tt} - a^2 \Delta u + b u_t = 0,
\end{align*}
其中 $a$ 为正常数，$b$ 为常数；

\item 横梁方程
\begin{align*}
u_t + u_{xxxx} = 0.
\end{align*}
\end{enumerate}
\end{examples}

\begin{examples}[较著名的一些非线性偏微分方程]
\begin{enumerate}[(1)]
\item 非线性 Poisson 方程
\begin{align*}
\Delta u = u^3 - u;
\end{align*}

\item 极小曲面方程
\begin{align*}
\operatorname{div}\left( \frac{D u}{(1 + |D u|^2)^{\frac{1}{2}}} \right) = 0;
\end{align*}

\item Monge-Amp\`{e}re 方程
\begin{align*}
\det(D^2 u) = f(x);
\end{align*}

\item Hamilton-Jacobi 方程
\begin{align*}
u_t + H(Du) = 0,
\end{align*}
其中 $H:\mathbb{R}^n\to\mathbb{R}$ 为已知函数；

\item Burgers 方程
\begin{align*}
u_t + u u_x = 0;
\end{align*}

\item 守恒方程
\begin{align*}
u_t + \operatorname{div} F(u) = 0;
\end{align*}

\item 多孔介质方程
\begin{align*}
u_t - \Delta u^\gamma = 0 \quad (\gamma > 1 \text{ 为常数});
\end{align*}

\item Korteweg-deVries (KdV) 方程
\begin{align*}
u_t + u u_x + u_{xxx} = 0;
\end{align*}

\item $p$-Laplace 方程
\begin{align*}
\operatorname{div} (|D u|^{p-2} D u) = 0 \quad (p > 1 \text{ 为常数});
\end{align*}

\item 非线性波动方程
\begin{align*}
u_{tt} - a^2 \Delta u = f(u) \quad (a > 0 \text{ 为常数});
\end{align*}

\item Boltzmann 方程
\begin{align*}
f_t + v \cdot D_x f = Q(f, f),
\end{align*}
其中 $f=f(x,v,t)$，$Q(f,f)=Q(f(x,v,t),f(x,v,t))$ 为碰撞项，这里
\begin{align*}
Q(\varphi,\varphi) = \int_{\mathbb{R}^n} \mathrm{d}v_* \int_{S^{n-1}} \mathrm{d}\omega [\varphi(v')\varphi(v'_*) - \varphi(v)\varphi(v_*)] B(v-v_*, \omega),
\end{align*}
$S^{n-1}$ 为 $n-1$ 维单位球面，$v'$, $v'_*$ 由等式 $v'=v-[(v-v_*)\cdot\omega]\omega$ 及 $v'_*=v+[(v-v_*)\cdot\omega]\omega$ 给定，$B$ 为碰撞核，$D_x f$ 表示 $f$ 对 $x$ 的偏导数。
\end{enumerate}
\end{examples}

\begin{examples}[较著名的一些线性偏微分方程组]
\begin{enumerate}[(1)]
\item 线性弹性平衡方程组
\begin{align*}
\mu \Delta u + (\lambda + \mu) D(\operatorname{div} u) = 0 \quad (\mu, \lambda > 0 \text{ 为常数});
\end{align*}

\item 线性弹性发展方程组
\begin{align*}
u_{tt} - \mu \Delta u - (\lambda + \mu) D(\operatorname{div} u) = 0 \quad (\mu, \lambda \text{ 为常数});
\end{align*}

\item Maxwell 方程组
\begin{align*}
\frac{1}{c} \cdot \frac{\partial E}{\partial t} = \operatorname{curl} B, \qquad
\frac{1}{c} \cdot \frac{\partial B}{\partial t} = -\operatorname{curl} E, \qquad
\operatorname{div} E = \operatorname{div} B = 0,
\end{align*}
这里 $E$, $B$ 分别为电场强度和磁场强度，$c$ 为光速。
\end{enumerate}
\end{examples}

\begin{examples}[较著名的一些非线性偏微分方程组]
\begin{enumerate}[(1)]
\item 守恒律方程组
\begin{align*}
u_t + [F(u)]_x = 0;
\end{align*}

\item 反应扩散方程组
\begin{align*}
u_t - a^2 \Delta u = f(u) \quad (a > 0 \text{ 为常数});
\end{align*}

\item Euler 方程组 (不可压无粘性流)
\begin{align*}
\begin{cases}
u_t + (u\cdot D)u + Dp = 0, \\
\operatorname{div} u = 0,
\end{cases}
\end{align*}
这里 $u$, $p$ 分别为流体的速度和压力；

\item Navier-Stokes 方程组 (不可压粘性流)
\begin{align*}
\begin{cases}
u_t + (u\cdot D)u - \mu \Delta u + Dp = 0, \\
\operatorname{div} u = 0,
\end{cases}
\end{align*}
其中 $\mu$ 为粘性系数，$u$, $p$ 分别为流体的速度和压力。
\end{enumerate}
\end{examples}

\begin{examples}[三个典型的二阶线性偏微分方程]
\begin{enumerate}[(1)]
\item 位势方程
\begin{align}
-\Delta u = f(x); \label{eq:potential}
\end{align}

\item 热方程
\begin{align}
u_t - a^2 \Delta u = f(x,t) \quad (a > 0 \text{ 为常数}); \label{eq:heat}
\end{align}

\item 波动方程
\begin{align}
u_{tt} - a^2 \Delta u = f(x,t) \quad (a > 0 \text{ 为常数}), \label{eq:wave}
\end{align}
其中 $\Delta = \sum_{i=1}^n \frac{\partial^2}{\partial x_i^2}$ 是 Laplace 算子。
\end{enumerate}
\end{examples}

\begin{definition}
在 $\mathbb{R}^m$ 中，设二阶线性偏微分方程的一般形式为
\begin{align}
\sum_{i,j=1}^{m} a_{ij}(x) u_{x_i x_j} + \sum_{i=1}^{m} b_i(x) u_{x_i} + c(x) u(x) = f(x),\label{eq:general-second1}
\end{align}
其中 $a_{ij}=a_{ji}$ $(i,j=1,\cdots,m)$,以 $A$ 表示矩阵 $(a_{ij})_{m\times m}$.

若矩阵 $A(x_0)$ 的 $m$ 个特征值都是负数，则称方程 \eqref{eq:general-second1} 在点 $x_0$ 属于\textbf{椭圆型}；

若矩阵 $A(x_0)$ 的 $m$ 个特征值除了一个特征值为 $0$ 外，其他 $m-1$ 个特征值都是负数，则称方程 \eqref{eq:general-second1} 在点 $x_0$ 属于\textbf{抛物型}；

若矩阵 $A(x_0)$ 的 $m$ 个特征值除了一个特征值为正数外，其他 $m-1$ 个特征值都是负数，则称方程 \eqref{eq:general-second1} 在点 $x_0$ 属于\textbf{双曲型}；

若对于区域 $\Omega$ 上的每一个点，方程 \eqref{eq:general-second1} 属于椭圆型，则称方程 \eqref{eq:general-second1} 在 $\Omega$ 上是\textbf{椭圆型的}；

若对于区域 $\Omega$ 上的每一个点，方程 \eqref{eq:general-second1} 属于抛物型，则称方程 \eqref{eq:general-second1} 在 $\Omega$ 上是\textbf{抛物型的}；

若对于区域 $\Omega$ 上的每一个点，方程 \eqref{eq:general-second1} 属于双曲型，则称方程 \eqref{eq:general-second1} 在 $\Omega$ 上是\textbf{双曲型的}。
\end{definition}
\begin{remark}
方程 \eqref{eq:potential}, \eqref{eq:heat} 和 \eqref{eq:wave} 是方程 \eqref{eq:general-second1}的特例。对于位势方程 \eqref{eq:potential}，取 $m=n$，则
\begin{align*}
A = \begin{pmatrix}
-1 & 0 & \cdots & 0 \\
0 & -1 & \cdots & 0 \\
\vdots & \vdots & \ddots & \vdots \\
0 & 0 & \cdots & -1
\end{pmatrix},
\end{align*}
且 $b_i(x)\equiv 0$ $(i=1,2,\cdots,n)$，$c(x)\equiv 0$；对于热方程 \eqref{eq:heat}，取 $m=n+1$，$t=x_{n+1}$，则
\begin{align*}
A = \begin{pmatrix}
-a^2 & \cdots & 0 & 0 \\
\vdots & \ddots & 0 & \vdots \\
0 & \cdots & -a^2 & 0 \\
0 & \cdots & 0 & 0
\end{pmatrix},
\end{align*}
且 $b_i(x)\equiv 0$ $(i=1,2,\cdots,n)$，$b_{n+1}(x)\equiv 1$，$c(x)\equiv 0$；对于波动方程 \eqref{eq:wave}，取 $m=n+1$，$t=x_{n+1}$，则
\begin{align*}
A = \begin{pmatrix}
-a^2 & \cdots & 0 & 0 \\
\vdots & \ddots & 0 & \vdots \\
0 & \cdots & -a^2 & 0 \\
0 & \cdots & 0 & 1
\end{pmatrix},
\end{align*}
且 $b_i(x)\equiv 0$ $(i=1,2,\cdots,n+1)$，$c(x)\equiv 0$。

从矩阵 $A$ 的特征值来考察方程 \eqref{eq:potential}, \eqref{eq:heat} 和 \eqref{eq:wave}，我们发现它们的差别在于：对于位势方程 \eqref{eq:potential}，系数矩阵 $A$ 的全部特征值都是负的，即 $A$ 是负定的；对于热方程 \eqref{eq:heat}，系数矩阵 $A$ 的全部特征值除了一个为 $0$ 外，其他都是负的，即 $A$ 是非正定的；对于波动方程 \eqref{eq:wave}，系数矩阵 $A$ 的全部特征值除了一个为正外，其他都是负的，即 $A$ 是不定的。

设 $x_0\in\mathbb{R}^m$，$A(x_0)$ 表示系数矩阵 $A$ 在点 $x_0$ 的值。由于系数矩阵 $A(x_0)$ 是对称矩阵，从线性代数的结论我们知道它总是可以对角化。基于此，我们按照上述定义对于一般二阶线性偏微分方程进行分类。
\end{remark}

当然，从矩阵 $A(x_0)$ 的性质分析还有其他的情形。由于这些方程的实际背景至今仍不清楚，远没有像上述三类方程引起人们的广泛兴趣。由上面的分类我们知道位势方程 \eqref{eq:potential} 是椭圆型方程，热方程 \eqref{eq:heat} 是抛物型方程，波动方程 \eqref{eq:wave} 是双曲型方程。方程 \eqref{eq:potential}, \eqref{eq:heat} 和 \eqref{eq:wave} 分别称为椭圆型、抛物型和双曲型方程的标准型。

\begin{theorem}
在 $\mathbb{R}^m$ 中，设二阶线性偏微分方程的为
\begin{align}
\sum_{i,j=1}^{m} a_{ij}(x) u_{x_i x_j} + \sum_{i=1}^{m} b_i(x) u_{x_i} + c(x) u(x) = f(x), \label{eq:general-second}
\end{align}
其中 $a_{ij}=a_{ji}$ $(i,j=1,\cdots,m)$,以 $A$ 表示矩阵 $(a_{ij})_{m\times m}$.
若方程 \eqref{eq:general-second} 的二阶项的系数矩阵 $A$ 是常数矩阵，且它属于椭圆型（或抛物型，双曲型）方程，则存在一个非奇异的自变量替换把方程 \eqref{eq:general-second} 的二阶项化为形如方程 \eqref{eq:potential}（或 \eqref{eq:heat}, \eqref{eq:wave}）的标准型。
\end{theorem}
\begin{proof}
我们只证明方程 \eqref{eq:general-second} 是椭圆型的情形，其他情形类似。将方程 \eqref{eq:general-second} 写成
\begin{align*}
\sum_{i,j=1}^{m} a_{ij} u_{x_i x_j}= f(x,u,D_x u),
\end{align*}
其中$D_x u$表示 $u$关于$x$的所有一阶项。由于方程 \eqref{eq:general-second} 属于椭圆型，因而一定存在一个非奇异矩阵 $B=(b_{kl})_{m\times m}$，使得系数矩阵 $A$ 可以对角化，且 $BAB^T = -I$（单位阵），即
\begin{align*}
\sum_{i,j=1}^{m} b_{ki} a_{ij} b_{lj} = -\delta_{kl}, \qquad k,l=1,\cdots,m,
\end{align*}
其中 $\delta_{kl}$ 是 Kronecker 符号。由于矩阵 $B$ 非奇异，作自变量替换
\begin{align*}
y_k = \sum_{i=1}^{m} b_{ki} x_i, \qquad k=1,\cdots,m,
\end{align*}
再由链式法则可将方程 \eqref{eq:general-second} 的二阶项化为
\begin{align*}
\sum_{i,j=1}^{m} a_{ij} u_{x_i x_j}
= \sum_{i,j,k,l=1}^{m} a_{ij} b_{ki} b_{lj} u_{y_k y_l}
= -\sum_{k,l=1}^{m} \delta_{kl} u_{y_k y_l}
= -\Delta_y u,
\end{align*}
从而方程 \eqref{eq:general-second} 对于新变量 $y_1,\cdots,y_m$ 具有如下形式：
\begin{align*}
-\Delta_y u = \tilde{f}(y,u,D_y u),
\end{align*}
这里 $D_y u$ 表示 $u$ 关于 $y$ 的所有一阶偏导数。定理至此获证。
\end{proof}









\end{document}